\section{What is the alternative to escalation?}
Suppose a 14B model answers every question and a 32B model replaces unusually long answers. On a held-out AIME archive, the gate scores 56.25\% under independent pairing. It gains 4.19 percentage points over escalating the same fraction of requests at random. Yet at equal prices per output token the gate spends more than calling the 32B model directly on every question, which scores 60.42\%. Both comparisons use exactly the same generations. They answer different questions for someone deciding how to spend an inference budget.

For budget allocation, a simple alternative chooses one direct service per request before reading the question. Randomizing between two services fills the line segment between their expected costs and scores. RouterBench uses this principle to benchmark routing \cite{routerbench}, and UniRoute fits a random direct comparator on validation data \cite{uniroute}. The additional difficulty here is that the cascade itself need not have a known position on that diagram. From standalone logs with repeated draws, changing which fallback response follows a long first answer moves the gate's cost and score together. Comparing one accuracy interval with another cost interval can price a replay that the records could never produce.

We study a fixed gate that pays for a completed first answer and can replace it with a fresh answer to the original prompt. Our main contribution is a sharp interval for its excess score over direct service \emph{at its own spend} when the response pairing is unknown. We distinguish a single draw-label convention for an entire benchmark from a different coupling on each question. A lower and upper computational certificate makes both the algebra and numerical answer checkable. The same construction yields price thresholds at which some or all alignments favor escalation. For independent calls the empirical product supplies a single comparison and needs no unknown-pairing assumption.

This problem sits between research on learned deferral from uncertainty features \cite{gupta}, stochastic coupling of model outputs for oracle routing \cite{chen1,chen2}, and selection-valid evaluation of deployable routing \cite{shihab}. Cascade routing optimizes routing and escalation strategies given quality and cost information \cite{cascade}. Here the policy is held fixed and its realized operating point is partially identified. We use classical assignment and minimax results \cite{kuhn,sion}. The application-specific result connects the feasible \emph{joint} cost and score of a post-answer gate to a budget-dependent comparator. We examine two repeated-generation archives and then apply the budget comparison to FrugalGPT and AutoMix \cite{frugal,automix}, where the authors have already measured policy costs and scores. These releases let us examine how often the choice changes, and by how much, even when there is no pairing problem to solve.
